% Generated. Do not edit.
% Projection scriptural-date-location/section, version 1: the Scriptural Date and
% Location page, generated whole from research/chronology.toml, the Scripture
% chronology corpus and the source library by tools/tpt document-types refresh.
% document: liturgy/roman-rite/1962/propers/temporal/58-eighteenth-after-pentecost
\runninghead{Scriptural Date and Location}
\section{Scriptural Date and Location}
\label{sec:date-location}

{\small Each passage the Mass reads or sings from Scripture, in canonical order. The Context column names what is dated --- the event the passage narrates, a setting it remembers, the traditional attribution of its author, or the composition of its book --- and where it belongs; the Date column gives the answer of the Scripture chronology corpus for that subject, with every alternative it records. Texts of the Mass that are not drawn from Scripture have no row. The evidence for each date follows the table.\par}
\medskip
% Generated. Do not edit.
% Deterministic projection of research/chronology.toml.
% schema: 2
% document: liturgy/roman-rite/1962/propers/temporal/58-eighteenth-after-pentecost
% generated-by: tools/tpt proper-chronology annotations
\newcommand{\chronologyannotationclaim}[6]{#6}
\newcommand{\chronologyannotationreach}[3]{}
\newcommand{\chronologyannotationgroup}[3]{#3}
\newcommand{\chronologyannotation}[1]{%
  \ifcsname triptychchronologyannotation@#1\endcsname
    \csname triptychchronologyannotation@#1\endcsname
  \else
    \PackageError{triptych}{No chronology annotation for #1}{}%
  \fi}
% introit: Introit
\expandafter\def\csname triptychchronologyannotation@introit\endcsname{%
  \chronologyannotationgroup{historical-setting}{research-pending}{\textbf{Historical setting} -- Occasion date unresolved.}%
  \space \chronologyannotationgroup{composition}{nonuniform}{\textbf{Composition} -- No single date applies across every cited locus.}%
}
% epistle: Epistle
\expandafter\def\csname triptychchronologyannotation@epistle\endcsname{%
  \chronologyannotationgroup{composition}{disputed}{\textbf{Composition} -- disputed: \chronologyannotationreach{1Cor.1.4}{1Cor}{true}\chronologyannotationreach{1Cor.1.5}{1Cor}{true}\chronologyannotationreach{1Cor.1.6}{1Cor}{true}\chronologyannotationreach{1Cor.1.7}{1Cor}{true}\chronologyannotationreach{1Cor.1.8}{1Cor}{true}\chronologyannotationclaim{composition.first-epistle-to-the-corinthians}{composition}{catholic-traditional-v1}{disputed}{(1 and 2 Corinthians; Galatians), 56}{A.D. 56}; \chronologyannotationreach{1Cor.1.4}{1Cor}{true}\chronologyannotationreach{1Cor.1.5}{1Cor}{true}\chronologyannotationreach{1Cor.1.6}{1Cor}{true}\chronologyannotationreach{1Cor.1.7}{1Cor}{true}\chronologyannotationreach{1Cor.1.8}{1Cor}{true}\chronologyannotationclaim{composition.first-epistle-to-the-corinthians}{composition}{catholic-traditional-v1}{disputed}{about Easter A.D. 57}{A.D. 57}.}%
}
% gradual: Gradual
\expandafter\def\csname triptychchronologyannotation@gradual\endcsname{%
  \chronologyannotationgroup{historical-setting}{research-pending}{\textbf{Historical setting} -- Occasion date unresolved.}%
  \space \chronologyannotationgroup{composition}{preferred}{\textbf{Composition}: \chronologyannotationreach{Ps.121.1}{Ps}{true}\chronologyannotationreach{Ps.121.7}{Ps}{true}\chronologyannotationclaim{critical.psalms.latest-composition-boundary}{composition}{catholic-critical-v1}{preferred}{before the Maccabean period, around 165 B.C.; no individual psalm can be dated securely}{Before c. 165 B.C.}}%
}
% alleluia: Alleluia
\expandafter\def\csname triptychchronologyannotation@alleluia\endcsname{%
  \chronologyannotationgroup{historical-setting}{research-pending}{\textbf{Historical setting} -- Occasion date unresolved.}%
  \space \chronologyannotationgroup{composition}{preferred}{\textbf{Composition}: \chronologyannotationreach{Ps.101.16}{Ps}{true}\chronologyannotationclaim{critical.psalms.latest-composition-boundary}{composition}{catholic-critical-v1}{preferred}{before the Maccabean period, around 165 B.C.; no individual psalm can be dated securely}{Before c. 165 B.C.}}%
}
% gospel: Gospel
\expandafter\def\csname triptychchronologyannotation@gospel\endcsname{%
  \chronologyannotationgroup{narrated-event}{research-pending}{\textbf{Event} -- No narrated-event date in the chronology corpus.}%
  \space \chronologyannotationgroup{composition}{disputed}{\textbf{Composition} -- disputed: \chronologyannotationreach{Matt.9.1}{Matt}{true}\chronologyannotationreach{Matt.9.2}{Matt}{true}\chronologyannotationreach{Matt.9.3}{Matt}{true}\chronologyannotationreach{Matt.9.4}{Matt}{true}\chronologyannotationreach{Matt.9.5}{Matt}{true}\chronologyannotationreach{Matt.9.6}{Matt}{true}\chronologyannotationreach{Matt.9.7}{Matt}{true}\chronologyannotationreach{Matt.9.8}{Matt}{true}\chronologyannotationclaim{composition.gospel-of-matthew}{composition}{catholic-traditional-v1}{disputed}{about A.D. 38-45}{c. A.D. 38--45}; \chronologyannotationreach{Matt.9.1}{Matt}{true}\chronologyannotationreach{Matt.9.2}{Matt}{true}\chronologyannotationreach{Matt.9.3}{Matt}{true}\chronologyannotationreach{Matt.9.4}{Matt}{true}\chronologyannotationreach{Matt.9.5}{Matt}{true}\chronologyannotationreach{Matt.9.6}{Matt}{true}\chronologyannotationreach{Matt.9.7}{Matt}{true}\chronologyannotationreach{Matt.9.8}{Matt}{true}\chronologyannotationclaim{composition.gospel-of-matthew}{composition}{catholic-traditional-v1}{disputed}{about the year 40-42}{c. A.D. 40--42}; \chronologyannotationreach{Matt.9.1}{Matt}{true}\chronologyannotationreach{Matt.9.2}{Matt}{true}\chronologyannotationreach{Matt.9.3}{Matt}{true}\chronologyannotationreach{Matt.9.4}{Matt}{true}\chronologyannotationreach{Matt.9.5}{Matt}{true}\chronologyannotationreach{Matt.9.6}{Matt}{true}\chronologyannotationreach{Matt.9.7}{Matt}{true}\chronologyannotationreach{Matt.9.8}{Matt}{true}\chronologyannotationclaim{composition.gospel-of-matthew}{composition}{catholic-traditional-v1}{disputed}{the years 40-45}{A.D. 40--45}; \chronologyannotationreach{Matt.9.1}{Matt}{true}\chronologyannotationreach{Matt.9.2}{Matt}{true}\chronologyannotationreach{Matt.9.3}{Matt}{true}\chronologyannotationreach{Matt.9.4}{Matt}{true}\chronologyannotationreach{Matt.9.5}{Matt}{true}\chronologyannotationreach{Matt.9.6}{Matt}{true}\chronologyannotationreach{Matt.9.7}{Matt}{true}\chronologyannotationreach{Matt.9.8}{Matt}{true}\chronologyannotationclaim{composition.gospel-of-matthew}{composition}{catholic-traditional-v1}{disputed}{about the year 60-68}{c. A.D. 60--68}; \chronologyannotationreach{Matt.9.1}{Matt}{true}\chronologyannotationreach{Matt.9.2}{Matt}{true}\chronologyannotationreach{Matt.9.3}{Matt}{true}\chronologyannotationreach{Matt.9.4}{Matt}{true}\chronologyannotationreach{Matt.9.5}{Matt}{true}\chronologyannotationreach{Matt.9.6}{Matt}{true}\chronologyannotationreach{Matt.9.7}{Matt}{true}\chronologyannotationreach{Matt.9.8}{Matt}{true}\chronologyannotationclaim{composition.gospel-of-matthew}{composition}{catholic-traditional-v1}{disputed}{about the years 64-67}{c. A.D. 64--67}; \chronologyannotationreach{Matt.9.1}{Matt}{true}\chronologyannotationreach{Matt.9.2}{Matt}{true}\chronologyannotationreach{Matt.9.3}{Matt}{true}\chronologyannotationreach{Matt.9.4}{Matt}{true}\chronologyannotationreach{Matt.9.5}{Matt}{true}\chronologyannotationreach{Matt.9.6}{Matt}{true}\chronologyannotationreach{Matt.9.7}{Matt}{true}\chronologyannotationreach{Matt.9.8}{Matt}{true}\chronologyannotationclaim{composition.gospel-of-matthew}{composition}{catholic-traditional-v1}{disputed}{about the year 50}{c. A.D. 50}.}%
}
% offertory: Offertory
\expandafter\def\csname triptychchronologyannotation@offertory\endcsname{%
  \chronologyannotationgroup{narrated-event}{preferred}{\textbf{Event}: \chronologyannotationreach{Ex.24.4}{Ex.24.1-11}{false}\chronologyannotationreach{Ex.24.5}{Ex.24.1-11}{false}\chronologyannotationclaim{israel.wilderness.covenant-at-sinai}{narrated-event}{catholic-traditional-v1}{preferred}{rising in the morning}{The sealing of the covenant at the foot of the mount, rising in the morning}.}%
}
% communion: Communion
\expandafter\def\csname triptychchronologyannotation@communion\endcsname{%
  \chronologyannotationgroup{superscription-setting}{preferred}{\textbf{Superscription setting}: \chronologyannotationreach{Ps.95.8}{Ps.95}{true}\chronologyannotationreach{Ps.95.9}{Ps.95}{true}\chronologyannotationclaim{israel.restoration.temple-raised-by-zorobabel}{superscription-setting}{catholic-traditional-v1}{preferred}{(537 B.C.)}{The raising of the second Temple by Zorobabel, B.C. 537}.}%
  \space \chronologyannotationgroup{composition}{preferred}{\textbf{Composition}: \chronologyannotationreach{Ps.95.8}{Ps}{true}\chronologyannotationreach{Ps.95.9}{Ps}{true}\chronologyannotationclaim{critical.psalms.latest-composition-boundary}{composition}{catholic-critical-v1}{preferred}{before the Maccabean period, around 165 B.C.; no individual psalm can be dated securely}{Before c. 165 B.C.}}%
  \space \chronologyannotationgroup{prophetic-referent}{disputed}{\textbf{Prophetic referent} -- disputed: \chronologyannotationreach{Ps.95.8}{Ps.67 Ps.95 Ps.96 Ps.97}{true}\chronologyannotationreach{Ps.95.9}{Ps.67 Ps.95 Ps.96 Ps.97}{true}\chronologyannotationclaim{life-of-christ.nativity}{prophetic-referent}{catholic-traditional-v1}{disputed}{the year of Rome 750 which he styles 3 B.C.}{The Nativity of Our Lord, B.C. 3}; \chronologyannotationreach{Ps.95.8}{Ps.67 Ps.95 Ps.96 Ps.97}{true}\chronologyannotationreach{Ps.95.9}{Ps.67 Ps.95 Ps.96 Ps.97}{true}\chronologyannotationclaim{life-of-christ.nativity}{prophetic-referent}{catholic-traditional-v1}{disputed}{probably the year 7 B.C.}{The Nativity of Our Lord, c. B.C. 7}.}%
}

% Generated. Do not edit.
% Complete orientation-v3 table; wording, dates and order belong to the corpus.
\begin{chronologysharedtable}
\chronologyorientationrow{row-51d2ec9a134cd457}{}
\midrule
\chronologyorientationrow{row-fcd5f131849b225b}{}
\chronologyorientationrow{row-bc7aef6490bda13d}{}
\chronologyorientationrow{row-91644ed46814e0f1}{}
\midrule
\chronologyorientationrow{row-25509eeb71cc905d}{}
\chronologyorientationrow{row-76ca94f9b8bd2f77}{}
\midrule
\chronologyorientationrow{row-50f2f83dc1dbafe8}{}
\chronologyorientationrow{row-8f040a8d598bacf5}{}
\midrule
\chronologyorientationrow{row-7a44ff28a5b1e6d3}{}
\chronologyorientationrow{row-bf3a7f963c760efa}{}
\midrule
\chronologyorientationrow{row-be28fc871b1307ae}{}
\midrule
\chronologyorientationrow{row-386afb378f8dcb34}{}
\chronologyorientationrow{row-f64244e956334aaa}{}
\midrule
\chronologyorientationrow{row-fcd145bad1cdb5b4}{}
\end{chronologysharedtable}
{\footnotesize Psalm numbering: Vg. = Vulgate; Heb. = Hebrew.\par}


\subsection{Evidence for each date}
\label{sec:date-location-evidence}
{\small For each passage, as the Scriptural Date and Location table names it: what is dated, how the corpus weighs each answer, the source's own words, and the sources cited. A subject the corpus does not date is reported as such and given no date.\par}
{\small
\par\smallskip\noindent\textit{Offertory: Ex. 24:4--5}\par\nobreak
\begin{itemize}[leftmargin=1.2em,itemsep=1pt,topsep=1pt,parsep=0pt]
\item \textbf{Narrated event} (preferred). The sealing of the covenant at the foot of the mount: Rising in the morning (The Holy Bible, Douay-Rheims (Challoner revision), Project Gutenberg eBook 1581: Ex. 24:3, 24:4).
\end{itemize}
\par\smallskip\noindent\textit{Communion: Ps. 95:8--9 (Vg.; Heb. 96)}\par\nobreak
\begin{itemize}[leftmargin=1.2em,itemsep=1pt,topsep=1pt,parsep=0pt]
\item \textbf{Superscription setting} (preferred). The raising of the second Temple by Zorobabel: B.C.~537 (``(537 B.C.)''; The Catholic Encyclopedia, Volume 14 (New York, 1912)). Place: Temple at Jerusalem (The Holy Bible, Douay-Rheims (Challoner revision), Project Gutenberg eBook 1581: Ps. 95:1, 1 Esd. 3:8).
\item \textbf{Composition} (preferred). Psalms: Before c.~165 B.C. (``before the Maccabean period, around 165 B.C.; no individual psalm can be dated securely''; Catholic critical chronology; New American Bible, Revised Edition --- USCCB web state (2026), Psalms introduction).
\item \textbf{Prophetic referent} (disputed). The Nativity of Our Lord: B.C.~3 (``the year of Rome 750 which he styles 3 B.C.''; The Catholic Encyclopedia, Volume 3 (New York, 1908)); c.~B.C.~7 (``probably the year 7 B.C.''; The Catholic Encyclopedia, Volume 3 (New York, 1908)). Place: Bethlehem as the Christian referent (The Holy Bible, Douay-Rheims (Challoner revision), Project Gutenberg eBook 1581: Matt. 2:1).
\end{itemize}
\par\smallskip\noindent\textit{Alleluia: Ps. 101:16 (Vg.; Heb. 102)}\par\nobreak
\begin{itemize}[leftmargin=1.2em,itemsep=1pt,topsep=1pt,parsep=0pt]
\item \textbf{Historical setting}. No date: the chronology corpus has not yet dated this (research pending).
\item \textbf{Composition}. As under Communion: Ps. 95:8--9 (Vg.; Heb. 96).
\end{itemize}
\par\smallskip\noindent\textit{Gradual: Ps. 121:1 (Vg.; Heb. 122), Ps. 121:7 (Vg.; Heb. 122)}\par\nobreak
\begin{itemize}[leftmargin=1.2em,itemsep=1pt,topsep=1pt,parsep=0pt]
\item \textbf{Historical setting}. As under Alleluia: Ps. 101:16 (Vg.; Heb. 102).
\item \textbf{Composition}. As under Communion: Ps. 95:8--9 (Vg.; Heb. 96).
\end{itemize}
\par\smallskip\noindent\textit{Introit: Ps. 121:1 (Vg.; Heb. 122)}\par\nobreak
\begin{itemize}[leftmargin=1.2em,itemsep=1pt,topsep=1pt,parsep=0pt]
\item \textbf{Historical setting}. As under Alleluia: Ps. 101:16 (Vg.; Heb. 102).
\item \textbf{Composition}. As under Communion: Ps. 95:8--9 (Vg.; Heb. 96).
\end{itemize}
\par\smallskip\noindent\textit{Introit: Ecclus. 36:18}\par\nobreak
\begin{itemize}[leftmargin=1.2em,itemsep=1pt,topsep=1pt,parsep=0pt]
\item \textbf{Composition} (answers differ). Ecclesiasticus: preferred: B.C.~190--170 (``between 190 and 170 B.C.''; The Catholic Encyclopedia, Volume 5 (New York, 1909)); alternative: c.~B.C.~280 (``about 280 B.C.''; The Catholic Encyclopedia, Volume 5 (New York, 1909)).
\end{itemize}
\par\smallskip\noindent\textit{Gospel: Matt. 9:1--8}\par\nobreak
\begin{itemize}[leftmargin=1.2em,itemsep=1pt,topsep=1pt,parsep=0pt]
\item \textbf{Narrated event}. No date: the chronology corpus has not yet dated this (research pending).
\item \textbf{Composition} (disputed). The Gospel of St Matthew: c.~A.D.~38--45 (``about A.D. 38-45''; The Catholic Encyclopedia, Volume 10 (1911); The Catholic Encyclopedia, Volume 14 (New York, 1912)); c.~A.D.~40--42 (``about the year 40-42''; The Catholic Encyclopedia, Volume 10 (1911)); A.D.~40--45 (``the years 40-45''; The Catholic Encyclopedia, Volume 10 (1911)); c.~A.D.~60--68 (``about the year 60-68''; The Catholic Encyclopedia, Volume 10 (1911)); c.~A.D.~64--67 (``about the years 64-67''; The Catholic Encyclopedia, Volume 10 (1911)); c.~A.D.~50 (``about the year 50''; The Catholic Encyclopedia, Volume 14 (New York, 1912)). Place: Semitic original associated with Palestine (The Catholic Encyclopedia, Volume 10 (1911)).
\end{itemize}
\par\smallskip\noindent\textit{Epistle: 1 Cor. 1:4--8}\par\nobreak
\begin{itemize}[leftmargin=1.2em,itemsep=1pt,topsep=1pt,parsep=0pt]
\item \textbf{Composition} (disputed). The First Epistle to the Corinthians: A.D.~56 (``(1 and 2 Corinthians; Galatians), 56''; The Catholic Encyclopedia, Volume 11 (New York, 1911)); A.D.~57 (``about Easter A.D. 57''; The Catholic Encyclopedia, Volume 4 (New York, 1908)).
\end{itemize}
\par}
